% ==============================================================
\section{結論}\label{sec:conclusion}
% ==============================================================

本研究では、Apriori-Window~\cite{aprioriwindow}を
1次元（時間軸）から3次元（時間 $\times$ 空間 $x \times$ 空間 $y$）に拡張した
\textbf{3次元 Apriori-Window アルゴリズム}を提案した。
3つの高速化技術---3次元接頭辞和、$t$ 軸ストライド最適化、Apriori 枝刈り---は
いずれも Apriori-Window の対応する1次元技術の直接拡張であり、
ナイーブ実装の $O(M \cdot G \cdot W_t W_x W_y)$ から
$O(M \cdot G + N \cdot W_x W_y)$ へ計算量を削減した。

合成データ実験では以下を確認した：
(1) \textbf{EX-A}: 3次元設定（$W_x=W_y=3$）は植え付け3パターンを
Precision $= 1.0$、Recall $= 1.0$ で完全回収したのに対し、
1次元相当設定（$W_x=W_y=50$）は背景ノイズ由来の偽陽性パターンを
多数検出した（Precision $\approx 0.03$）。
空間次元の明示的なモデル化が精度に本質的に重要であることを示す。
(2) \textbf{EX-B}: $N = 100{,}000$ トランザクションで約 1.4 秒を達成し、
$N$ に対してほぼ線形のスケーリングを示した。
(3) \textbf{EX-C}: Apriori 枝刈りにより $k \geq 3$ の不要な計算が
完全に排除された。
(4) \textbf{EX-D}: 比較実験（5手法、5シード）では、
3次元 AW のみが Precision $= 1.0$、Recall $= 1.0$（F1 $= 1.0$）を
全シードで達成した。
FIM\_global（Precision $= 0.016$）は空間局在化不可、
1D\_AW（F1平均 $= 0.971$）は時刻近接による偽陽性が発生、
Span\_3D（F1平均 $= 0.840$）は点密度の厳格さによる見逃しが生じた。

\paragraph{制約と限界.}
(1) 空間次元はグリッド離散化を前提とし、ビン幅の選択が結果に影響する。
(2) 空間ウィンドウサイズが大きい場合（$W_x W_y \gtrsim 100$）、
XY イベントリスト構築コストが支配的となり実行時間が増大する。
(3) 本手法は密集領域の検出に特化しており、
検出された密集領域と外部イベントとの統計的関連付けは行わない。

\paragraph{今後の課題.}
(1) 実データ（小売 POS データ、疫学データ等）への応用と評価。
(2) 検出された密集領域を外部イベントに帰属させるパイプライン
（Apriori-Window~\cite{aprioriwindow}の後続研究~\cite{dsaa2025}）への接続。
(3) ストリーミング設定へのオンライン更新対応。
(4) 適応的グリッド解像度選択。
